

%% Basierend auf einer TeXnicCenter-Vorlage von Mark Müller
%%%%%%%%%%%%%%%%%%%%%%%%%%%%%%%%%%%%%%%%%%%%%%%%%%%%%%%%%%%%%%%%%%%%%%%

% Wählen Sie die Optionen aus, indem Sie % vor der Option entfernen  
% Dokumentation des KOMA-Script-Packets: scrguide

%%%%%%%%%%%%%%%%%%%%%%%%%%%%%%%%%%%%%%%%%%%%%%%%%%%%%%%%%%%%%%%%%%%%%%%
%% Optionen zum Layout des Artikels                                  %%
%%%%%%%%%%%%%%%%%%%%%%%%%%%%%%%%%%%%%%%%%%%%%%%%%%%%%%%%%%%%%%%%%%%%%%%
\documentclass[%
%a5paper,							% alle weiteren Papierformat einstellbar
%landscape,						% Querformat
12pt,								% Schriftgröße (12pt, 11pt (Standard))
%BCOR1cm,							% Bindekorrektur, bspw. 1 cm
%DIVcalc,							% führt die Satzspiegelberechnung neu aus
%											  s. scrguide 2.4
%twoside,							% Doppelseiten
%twocolumn,						% zweispaltiger Satz
%halfparskip*,				% Absatzformatierung s. scrguide 3.1
%headsepline,					% Trennline zum Seitenkopf	
%footsepline,					% Trennline zum Seitenfuß
titlepage,						% Titelei auf eigener Seite
%normalheadings,			% Überschriften etwas kleiner (smallheadings)
%idxtotoc,						% Index im Inhaltsverzeichnis
%liststotoc,					% Abb.- und Tab.verzeichnis im Inhalt
%bibtotoc,						% Literaturverzeichnis im Inhalt
%abstracton,					% Überschrift über der Zusammenfassung an	
%leqno,   						% Nummerierung von Gleichungen links
%fleqn,								% Ausgabe von Gleichungen linksbündig
%draft								% überlangen Zeilen in Ausgabe gekennzeichnet
bibliography=totoc
]
{scrartcl}

\usepackage{geometry}
\geometry{a4paper, left=45mm, right=20mm, top=30mm, bottom=30mm}
% http://www.macuser.de/forum/f19/latex-zeilenabstand-aendern-189206/
\usepackage{setspace}

% Hurenkinder und Schusterjungen verhindern
\clubpenalty10000
\widowpenalty10000
\displaywidowpenalty=10000

\usepackage[flushmargin,hang,ragged]{footmisc}

%\pagestyle{empty}		% keine Kopf und Fußzeile (k. Seitenzahl)
%\pagestyle{headings}	% lebender Kolumnentitel  


%% Deutsche Anpassungen %%%%%%%%%%%%%%%%%%%%%%%%%%%%%%%%%%%%%
\usepackage[ngerman]{babel}
\usepackage[T1]{fontenc}
\usepackage[ansinew]{inputenc}
\usepackage{graphicx}

\usepackage{lmodern} %Type1-Schriftart für nicht-englische Texte
\usepackage{fancyhdr}

%% Für Code-Highlighting
%% http://www.ctan.org/tex-archive/macros/latex/contrib/minted/
\usepackage{minted}
\definecolor{bg}{rgb}{0.95,0.95,0.95}

%\usepackage{cite}
%\renewcommand{\citeleft}{}
%\renewcommand{\citeright}{}

%% Packages für Grafiken & Abbildungen %%%%%%%%%%%%%%%%%%%%%%
\usepackage{graphicx} %%Zum Laden von Grafiken
%\usepackage{subfig} %%Teilabbildungen in einer Abbildung
%\usepackage{pst-all} %%PSTricks - nicht verwendbar mit pdfLaTeX

%% Beachten Sie:
%% Die Einbindung einer Grafik erfolgt mit \includegraphics{Dateiname}
%% bzw. über den Dialog im Einfügen-Menü.
%% 
%% Im Modus "LaTeX => PDF" können Sie u.a. folgende Grafikformate verwenden:
%%   .jpg  .png  .pdf  .mps
%% 
%% In den Modi "LaTeX => DVI", "LaTeX => PS" und "LaTeX => PS => PDF"
%% können Sie u.a. folgende Grafikformate verwenden:
%%   .eps  .ps  .bmp  .pict  .pntg

%% Bibliographiestil %%%%%%%%%%%%%%%%%%%%%%%%%%%%%%%%%%%%%%%%%%%%%%%%%%
\usepackage[style=authoryear,backend=bibtex8]{biblatex}
\usepackage[babel,german=quotes]{csquotes}
%\addbibresource{Literatur.bib}
\bibliography{./Literatur}

\begin{document}
\parskip=1em
\parindent=0cm
\pagestyle{empty} %%Keine Kopf-/Fusszeilen auf den ersten Seiten.


%%%%%%%%%%%%%%%%%%%%%%%%%%%%%%%%%%%%%%%%%%%%%%%%%%%%%%%%%%%%%%%%%%%%%%%
%% Ihr Artikel                                                       %%
%%%%%%%%%%%%%%%%%%%%%%%%%%%%%%%%%%%%%%%%%%%%%%%%%%%%%%%%%%%%%%%%%%%%%%%

%% eigene Titelseitengestaltung %%%%%%%%%%%%%%%%%%%%%%%%%%%%%%%%%%%%%%%    
%\begin{titlepage}
%Einsetzen der TXC Vorlage "Deckblatt" möglich
%description: Deckblatt in Deutsch
%% Basierend auf einer TeXnicCenter-Vorlage von Tino Weinkauf.
%%%%%%%%%%%%%%%%%%%%%%%%%%%%%%%%%%%%%%%%%%%%%%%%%%%%%%%%%%%%%%

%%%%%%%%%%%%%%%%%%%%%%%%%%%%%%%%%%%%%%%%%%%%%%%%%%%%%%%%%%%%%
%% Deckblatt
%%%%%%%%%%%%%%%%%%%%%%%%%%%%%%%%%%%%%%%%%%%%%%%%%%%%%%%%%%%%%
%%
%% ACHTUNG: Sie ben�tigen ein Hauptdokument, um diese Datei
%%          benutzen zu k�nnen. Verwenden Sie im Hauptdokument
%%          den Befehl "\input{dateiname}", um diese
%%          Datei einzubinden.
%%

\begin{titlepage}
\vspace{5cm}
Stefan Waidele \hfill{} Matthias Vongerichten\\
Immatrikulationsnummer 1028171 \hfill{} Immatrikulationsnummer 23422001\\
Ensisheimer Stra�e 2 \hfill{} Deine Stra�e 5\\
79395 Neuenburg am Rhein \hfill{} 12345 Wherever\\
Stefan.Waidele@AKAD.de \hfill{} Matthias.Vongerichten@AKAD.de

\vfill

\begin{tabbing}
Modul INT02 --- \=Einf�hrung in die Internetprogrammierung\\ \\
                \>Assignment  
\end{tabbing}
\LARGE
\textsc{Erstellung einer Website \\f�r ein fiktives Busunternehmen}\\

\vfill

\normalsize
Betreuer: Andr� Langbein

\today %%Datum der Abgabe - am besten selbst reinschreiben.

\vfill

\includegraphics[width=3cm]{AKAD-Logo}

AKAD Hochschule Stuttgart

\end{titlepage}
%\end{titlepage}

%% Erzeugung von Verzeichnissen %%%%%%%%%%%%%%%%%%%%%%%%%%%%%%%%%%%%%%%
\setcounter{tocdepth}{2}
\tableofcontents			% Inhaltsverzeichnis
%\listoftables				% Tabellenverzeichnis
\listoffigures				% Abbildungsverzeichnis
\clearpage


%% Der Text %%%%%%%%%%%%%%%%%%%%%%%%%%%%%%%%%%%%%%%%%%%%%%%%%%%%%%%%%%%
\pagestyle{headings}
\begin{onehalfspacing}
\section{Einleitung}
\subsection{Ziel der Arbeit}

Ziel dieser Arbeit ist ...

\subsection{Vorgehensweise}


\begin{itemize}
	\item Grundlagen werden in den Konzepten kurz beschrieben.
	\item Umsetzung mit PHP, HTML, CSS, Apache Webserver
\end{itemize}


\subsection{Abgrenzung}

Nicht alle Seiten des Internetauftritts...

\begin{itemize}
	\item \textbf{Redaktionelles Konzept:}\\
	...
	\item \textbf{Navigationskonzept:}\\
	...
	\item \textbf{Gestaltungskonzept:}\\
	...

\end{itemize}

Es wird kein Social--Media--Konzept erstellt, da dies den Rahmen der Arbeit überschreiten würde. Allerdings wird die Einbindung von Web 2.0 Komponenten im Rahmen des Redaktionellen Konzepts betrachtet.
\section{Anforderungsanalyse}

\textsc{Stefan Waidele}

\subsection{Ziele der Website}

Die Unternehmenswebsite f�gt sich in die gro�e Bandbreite der Werbema�nahmen eines Unternehmens ein.
Daher muss diese sich auch nach den Grunds�tzen der Werbung und des Marketings richten.

Die allgemeinen Ziele der Website sind die Gewinnung neuer Kunden sowie die Erh�hung des Bekantheitsgrades des Unternhemens.

F�r diese Arbeit gehen wir von folgendenden Vorgaben\footnote{\cite{kloss}, S. 183ff} aus:

\begin{itemize}
	\item \textbf{Zielgruppe:}\\
	Das Unternehmen betrachtet Alleinreisende und Paare im Alter zwischen 40 und 65 Jahren als die Hauptzielgruppe f�r die angebotenen St�dtereisen. Au�erdem sollen Vereine und Reisegruppen gemischten Alters angesprochen werden. 
	Altersunabh�ngig ist die Zielgruppe haupts�chlich den Sinus-Charakteristiken "`B�rgerliche Mitte"', "`Traditionelle"' oder "`Konservativ--etablierte"' zuzuordnen\footnote{\cite{sinus}}.
	
	\item \textbf{Phase der Buchungsentscheidung:}\\
	Die Website soll dem Kunden zu jedem Zeitpunkt des Entscheidungsvorgangs\footnote{vgl. \cite{freyer}, S105} etwas bieten k�nnen. Es ist jedoch besonders auf die Informations-- und Handlungsphase zu achten, da hier der Einfluss des Unternehmens am direktesten ist.
\end{itemize}

\subsection{Funktionale Anforderungen}

Folgende Funktionen sollten dem Besucher der Seite angeboten werden:

\begin{itemize}
	\item Vermitteln von Informationen �ber das Unternehmen
	\item Darstellung der Reiseziele und deren Details
	\item Buchung von Reisen
	\item Kontaktaufnahme mit dem Unternehmen
\end{itemize}

\subsection{Nicht-funktionale Anforderungen}

Folgenden nicht--funktionalen bzw.qualitativen Kriterien soll die Website entsprechen: 

Abgestimmtes Erscheinungsbild, kurze Ladezeiten, 
einfaches Zurechtfinden, einfache Kontaktaufnahme, wenig Administrationsaufwand, mehrsprachigkeit (Deutsch, Englisch)

Dar�ber hinaus sind die Dialoggrunds�tze f�r interaktive Systeme gem�� ISO 9241-110 zu beachten\footnote{\cite{bhofmann}}:

Aufgabenangemessenheit, Selbstbeschreibungsf�higkeit, Erwartungskonformit�t, 
Fehlertoleranz, Steuerbarkeit, Individualisierbarkeit, Lernf�rderlichkeit

\subsection{Technische Voraussetzungen}

Zum Betreiben der Website wird ein Webspace mit mindestens 50MB Speicherplatz, vorzugsweise unlimitiertem Traffic, FTP-Zugang, und der M�glichkeit der mod\_rewrite--Aktivierung ben�tigt. Au�erdem muss der verwendete Webserver PHP interpretieren k�nnen. Eine Datenbank ist nicht erforderlich. F�r das Bearbeiten der Seiteninhalte gen�gt ein Text--Editor.
Zur besseren Auffindbarkeit der Website wird eine Top-Level-Domain empfohlen. Diese kann dem Webspace oft einfach zugebucht werden.
\section{Redaktionelles Konzept}

\textsc{Matthias Vongerichten}

\subsection{Start}

\textsc{Updatezyklus: j�hrlich}

Die Startseite hei�t den Besucher auf der Website willkommen. Hierf�r wird ein mehrzeiliger Flie�text angezeigt. In diesem wird die Vision und die Leidenschaft der Unternehmung beschrieben. Zus�tzlich wird im Flie�text ein Foto einer k�rzlich durchgef�hrten Reise angezeigt.

\subsection{�ber uns}

\textsc{Updatezyklus: bei Notwendigkeit}

Diese Seite wird auf 3 Abschnitte unterteilt: 

\begin{itemize}
\item \textbf{Abschnitt: Was uns auszeichnet}\\
Gefordert ist hier, dass herauskristallisiert wird was einem von der Konkurrenz unterscheidet bzw. abhebt. Auf was wird besonderen Wert gelegt? F�r was ist man bekannt?
\item \textbf{Abschnitt: Das Team ...}\\
Hier wird das Team in Form eines Gruppenbilds vorgestellt.
\item \textbf{Abschnitt: ... und unsere Geschichte}\\
Historische Entwicklung des Unternehmens. Was waren die "`Meilensteine"'? Wie kam es zur Gr�ndung? Auf was hatte man sich am Anfang spezialisiert und welchen Umfang bietet man heute? 
\end{itemize}

\subsection{Reiseziele}

\textsc{Updatezyklus: bei Notwendigkeit}

Hier werden die momentan besonders hervorzuhebenden Reisen, sortiert nach Zielen in Deutschland und Europa aufgelistet.
Die angebotenen St�dtereisen werden als kleiner Ausschnitt dargestellt. Der Ausschnitt zeigt ein aussagekr�ftiges Bild, eine kurze Beschreibung sowie einen Hyperlink zu den Reisedetails.

Auf Unterseiten werden nach L�ndern gegliedert alle angebotenen Reisen angezeigt. 

\subsection{Reisedetails}

\textsc{Updatezyklus: bei Notwendigkeit}

Hier erh�lt man genauere Informationen �ber eine St�dtereise. 
Auf der linken Seite sind mindestens 2 kleine Bilder vorhanden, welche vom Beschreibungstext umflossen werden. Der Text besteht aus einer kurzen Einf�hrung am Anfang und geht in die Tagesplanung �ber. In der Sidebar auf der rechten Seite erfolgt eine detaillierte Auflistung der Reisedaten. Folgende Elemente \footnote{\cite{freyer}, S103}, die die Entscheidungsphase ma�geblich unterst�tzen, sollen auf der Seite Erw�hnung finden: 

Reisezeitpunkt, Reisedauer, Reisegebiet, Reisegestaltung (Hier handelt es sich immer um Pauschalangebote), Kosten der Reise pro Person, Reisekomfort (Angabe der verwendeten Busklasse, Unterkunftsart, Nebenleistungen und -ausgaben.

Zus�tzlich wird ein Hyperlink zur Buchungsanfrage angezeigt. Der Link verweist hauf das �bliche Kontaktformular und �bergibt hierbei noch die ausgew�hlte St�dtereise als Parameter.

\clearpage
\subsection{Die Busse}

\textsc{Updatezyklus: bei Notwendigkeit}

Hier kann sich der potenzielle Kunde einen �berblick �ber den Fuhrpark des Unternehmens verschaffen. Ein Bild je Bus sowie dessen besonderen Merkmale stellen ein Minimum dar.

\subsection{Kontakt}

\textsc{Updatezyklus: bei Notwendigkeit}

Die Seite erm�glicht es dem Anwender auf verschiedene Weise mit dem Unternehmen in Kontakt zu treten. Angeboten wird:

Ein Kontaktformular, die E--Mail Adresse, Telefonnummer, Faxnummer, Links zu Social Media Plattformen

\subsection{Impressum}

\textsc{Updatezyklus: bei Notwendigkeit}

Nach Paragraf 5 des Telemediengesetz\footnote{\cite{justlaw}} sind folgende Angaben notwendig:

Name und Anschrift, E--Mail Adresse, Telefon, ggf. Register und Registernummer, 
ggf. Umsatzsteuer-- oder Wirtschaftssteueridentifikationsnummer

\subsection{AGB}

\textsc{Updatezyklus: bei Notwendigkeit}

Allgemeine Gesch�ftsbedingungen erweitern gesetzliche Regelungen nach eigenem Ermessen (innerhalb "`der wesentlichen Grundgedanken"') und sind daher freiwillig anzugeben. Im Falle dieses Busunternehmens k�nnte grunds�tzlich auf eine AGB verzichtet werden, jedoch bietet diese Seite auch die M�glichkeit sonstigen Informationspflichten bei Internetangeboten nachzugehen\footnote{\cite{tschwenke}}:

Widerrufsbelehrungspflichten gegen�ber Verbrauchern, Hinweise wie der Vertrag zustande kommt, 
Informationspflichten betreffend Preisangaben, Datenschutzhinweise, Hinweise nach der 
Dienstleistungs-Informationspflichten-Verordnung

\subsection{Social Media}

Eine Pr�senz in den sozialen Netzwerken geh�rt heutzutage in jedes erfolgreiche Website-Konzept. Neben der einfachen M�glichkeit von Kunden offiziell "`empfohlen"' zu werden bietet sich die M�glichkeit, als Alternative zu einem Weblog, potenzielle Kunden mit Neuigkeiten zu versorgen und mit Ihnen schnell in Kontakt zu treten. Aufgrund seiner aktuell noch dominierenden Stellung wird ein Account bei Facebook empfohlen - bei steigender Nutzeranzahl sp�ter auch ein Google+-Account, welcher dann mit Facebook integriert werden kann.

Im Header der Site kann die Facebook-Page besucht und "`geliked"' sowie bei Google+ - da dort noch kein eigener Account existiert - die Domain "`+1"'-ed werden. 
\section{Navigationskonzept}

\textsc{Matthias Vongerichten}

\subsection{Technische Umsetzung einheitlicher Navigation}

F�r jede Seite existiert eine PHP-Datei (Content) mit dem jeweiligen Inhalt, die au�erdem den Header und den Footer einbindet. Header und Footer werden somit f�r alle Seiten vereinheitlicht dargestellt und k�nnen zentral bearbeitet werden. 

\subsection{Seitenhierarchie}

Die Website besteht aus max. 2 Ebenen. Lediglich die Reisedetails geh�ren der zweiten Ebene an. Auf diese gelangt man �ber die Reiseziele.

\subsection{Hauptmen�}

Das Hauptmen� ist im Header angesiedelt und beinhaltet ausschlie�lich Hyperlinks zu Seiten der ersten Ebene.

\subsection{Breadcrumb Navigation}

Zur besseren Orientierung wird �ber dem Titel einer Seite eine "`Brotkr�mel"'-Navigation angezeigt. Diese ist im Content angesiedelt.

\subsection{Verwandte Seiten}

Je Reiseziel werden in einem gesonderten Abschnitt Hyperlinks zu relevanten Internetseiten angezeigt. Diese beziehen sich immer auf die ausgew�hlte Reise und werden daher im Content realisiert.

\subsection{Suchmaschinenoptimierung - SEO}

Um die Popularit�t der Site zu erh�hen werden heutzutage kaum noch META-tags verwendet\footnote{vgl. \cite{vertommen}}. Sinnvoll sind diese nur noch als Vorgabe des Beschreibungstextes in den Suchergebnissen  ("`description"'-tag) und bei einer gew�nschten Einschr�nkung der Suchmaschinen-Robots ("`robots"'-tag). Da eine Einschr�nkung der Suchmaschinen nicht erw�nscht ist, sollte lediglich das "`description"'-tag verwendet werden.

Um die Relevanz der Site zu erh�hen\footnote{vgl. \cite{googseo}}, sollte(n)

\begin{itemize}
	\item der Inhalt regelm��ig aktualisiert werden (siehe redaktionelles Konzept)
	\item die Reisen ausf�hrlich beschrieben sein
	\item die URL gut lesbar sein (Einsatz von mod\_rewrite)
	\item daf�r gesorgt werden, dass die Website auf anderen, qualitativ hochwertigen Seiten verlinkt wurde
	\item es die M�glichkeit geben die Site zu teilen (Social Media: Facebook, Google+ usw.)
\end{itemize}

Der Einsatz von mod\_rewrite erm�glicht es URL's so zu ver�ndern, dass diese besser lesbar und somit auch einpr�gsamer werden. Hierzu sind einige Modifikationen in der .htaccess-Datei notwendig.

\subsubsection{Sitemap}

Eine hierarchische Auflistung aller Einzeldokumente der Website erleichtert den Suchmaschinen das Indexieren und bietet den Benutzern au�erdem einen schnellen �berblick.




\section{Designkonzept}

\textsc{Matthias Vongerichten bzw. Stefan Waidele}

\subsection{Einheitliche Darstellung durch CSS}

...

\subsection{Gliederung der Seiten durch DIVs}

\subsubsection{Beurteilung TABLE}

Layout mit HTML-Tabellen ist B�SE!

\subsubsection{Beurteilung FRAME}

Layout mit HTML-Frames ist SCHLECHT!

\subsubsection{Beurteilung CSS-DIV}

Layout mit DIVs per CSS ist GUT!

(Sematisches Markup, Reihenfolge weniger relevant, erm�glicht Liquid/Responsive Layout, ...

\subsection{Farbkonzept}

...

\subsection{Richtlinien f�r Grafiken}

ALT-Tags

\subsection{Responsive Layout}

Mit Breakpoints f�r 1920px (HD), 960px (Klasisch Desktop), 800px (Beamer, Tablet), 480px (Smartphone)

Zwischenstufen Liquid.

\subsection{Special Effects}

Animationen (Dezent, evt. "`umdrehen"' bei Men�-MouseOver, Dezentes leuchten der Buchungs-Telefonnummer,...)

\subsection{Browserkompatibilit�t}

How far back in time do we want to travell?

\subsection{Noch Konzept...}

...

%%\section{Beispielseite: Startseite - MV}

\textsc{Mathias Vongerichten bzw. Stefan Waidele}

\subsection{�berschrift}

Erl�uterungen zur redaktionellen, gestalterischen und technischen Umsetzung.


%%\section{Beispielseite: L�nder�bersicht - MV}

\textsc{Matthias Vongerichten bzw. Stefan Waidele}

\subsection{�berschrift}

Erl�uterungen zur redaktionellen, gestalterischen und technischen Umsetzung.

Die von uns erstellten Seiten sollten hier im Detail mit Screenshots und Alternativen vorgestellt werden. Seiten, die wir nicht erstellen k�nnen wir in einem Kapitel zusammenfassen.
\section{Beispielseite: Reiseziele}
\label{bspreisen}
\textsc{Stefan Waidele}

\begin{figure}[h]
\includegraphics[width=\hsize]{scr_reisen}
\caption{Screenshot: �bersichtsseite Reisen}
\label{scr:reisen}
\end{figure}

Die Seite "`Reiseziele"' listet eine Auswahl der beliebtesten Reiseziele in Deutschland und Europa auf. 
Die nach L�ndern sortierten �bersichtsseiten sind hier in der Seitenleiste verlinkt. Die Seite wird 
von der Datei \code{/reiseziele/index.php} erzeugt.

Die hier gezeigten Elemente, besonders das Einbinden der gemeinsamen Elemente in \code{header.php} sowie \code{footer.php}
aber auch der individuelle Content--Bereich und die Seitenleite sind bei allen Seiten dieses Projekts gleich.

\subsection{Aufruf von header.php}

\begin{figure}[h]
\begin{minted}[bgcolor=bg]{php}
<?php 
$title="Reiseziele";
$desc = "�bersicht �ber unsere Reiseziele - Deutschland - Europa";
include($_SERVER['DOCUMENT_ROOT'].'/header.php'); 
?>
\end{minted}
\caption{Quellcode: Aufruf von header.php (PHP)}
\label{abb:header}
\end{figure}

Zun�chst werden hierin einige Angaben wie der Seitentitel in PHP-Variabeln abgelegt, um
dann die f�r alle Seiten gleiche Header-Datei \code{/header.php} aufzurufen (Vergleiche Abbildung~\ref{abb:header}). Hierin ist der gesamte HTML-Head Bereich 
sowie der f�r Besucher sichtbare Seitenkopf bis zum Navigationsmen� gespeichert. Der Ersteller oder Bearbeiter
der Reise�bersichtsseite braucht sich hierum nicht zu k�mmern.

\subsection{Realisierung der Breadcrumb--Navigation}

Nachdem die beiden f�r das Layout notwendigen DIVs ge�ffnet wurden, muss nun die Breadcrumb--Navigation angezeigt werden, anhand derer der Besucher schnell sehen kann, an welcher Stelle der Hierarchiestufe er sich befindet. Die zuvor besuchten, h�her liegenden Seiten sollen verlinkt sein. Der HTML--Code ist wie in Abbildung~\ref{abb:breadcrumb} zu gestalten. 

\begin{figure}[h]
\begin{minted}[bgcolor=bg]{html}
<div id="breadcrump">
   <a href="/" title="Zur�ck zur Startseite">Start</a> 
   > <a href="/reiseziele/" title="Kurz�bersicht">Reiseziele</a> 
   > Reisen in Deutschland
</div>
\end{minted}
\caption{Quellcode: Breadcrumb--Navigation (HTML)}
\label{abb:breadcrumb}
\end{figure}

\subsection{Kurzvorstellung der Reiseziele}

Nachdem der Seitentitel mit dem HTML-Tag \code{<h1>} und die Reiseregion (Deutschland bzw. Europa) mit \code{<h2>} ausgezeichnet wurden, beginnt die Beschreibung der eigentlichen Reiseziele mit einer �berschrift der Stufe \code{<h3>}, gefolgt von einem passenden Bild und dem beschreibenden Text. Am Ende der Kurzbeschreibung wird auf die Seite mit den entsprechenden Detailinformationen verlinkt. Somit ergibt sich eine Abfolge wie in  Abbildung~\ref{abb:stadt} zu sehen ist:

\begin{figure}[h]
\begin{minted}[bgcolor=bg]{html}
<div class="excerpt clearfix">
   <a href=""><h3>Berlin</h3></a>
   <a href=""><img src="/img/berlin1.jpg" 
                   alt="Fernsehturm in Berlin" /></a>
   <p>Lorem ipsum dolor sit amet, ...</p>
   <p><a href="">Weitere Informationen...</a></p>
</div>
\end{minted}
\caption{Quellcode: Reiseinformationen (HTML)}
\label{abb:stadt}
\end{figure}

Hierbei ist zu beachten, dass die Detailseite nicht nur am Ende des Textes, sondern auch bei der �berschrift und beim Bild verlinkt ist, um den Kunden m�glichst direkt zu den gew�nschten Informationen zu bringen.

Dieses Muster kann dann f�r jedes darzustellende Reiseziel wiederholt werden.
Aufgrund der Formatierungen im CSS muss der Redakteur wohl den korrekten HTML-Syntax sicherstellen, es befinden sich jedoch kaum Layout-Anweisungen im HTML-Code. Somit wird ein einheitliches Erscheinungsbild erzielt.

\subsection{Seitenleiste}

Auch die Seitenleiste wird in den entsprechenden Layout-DIVs nach einem einfachen, sich wiederholenden Muster realisiert.
Wie in Abbildung~\ref{abb:sidebar} zu erkennen ist, bestehen die �berschriften der Seitenleisten aus einem als \code{DIV} ausgezeichnetem Text, der von den Unterpunkten in Form einer HTML-Liste gefolgt wird. 

\begin{figure}[h]
\begin{minted}[bgcolor=bg]{html}
<div id="sidebar" class="clearfix">
   <div class="sidebar-topic">Reiseziele
   <span class="corner-right"></span></div>
      <ul>
         <li><a href="">Reisen in Deutschland</a></li>
         <li><a href="">L�nder�bersicht Europa</a></li>
      </ul>
   <div class="sidebar-topic">Bewertungen
   <span class="corner-right"></span></div>
      <ul>
         <li>5 Sterne in der Kundenbefragung</li>
         <li>7 Sonnen bei Holliday-Testsheck</li>
      </ul>
</div>
\end{minted}
\caption{Quellcode: Seitenleiste (HTML)}
\label{abb:sidebar}
\end{figure}

\subsection{Aufruf von footer.php}

Zum Abschluss der Seite wird mit der Codezeile aus Abbildung~\ref{abb:footer} noch die Datei \code{/footer.php} eingelesen. Hier wird der einheitliche Seitenfu� mit den Links zum Impressum, zu den AGBs und zur Sitemap dargestellt. 

\begin{figure}[h]
\begin{minted}[bgcolor=bg]{php}
<?php include($_SERVER['DOCUMENT_ROOT'].'/footer.php'); ?>
\end{minted}
\caption{Quellcode: Aufruf von footer.php (PHP)}
\label{abb:footer}
\end{figure}
\clearpage

\section{Beispielseite: Kontakt}

\textsc{Matthias Vongerichten}

%\includegraphics{scr_kontakt}
%Bitte Bild einf�gen, k.A. wie das geht :P

Die Seite "`Kontakt"' besteht zum gr��ten Teil aus einem Kontaktformular und den alternativen Kontaktm�glichkeiten in der Sidebar. 

Das Kontaktformular verwendet das �bliche HTML-Tag "`Form"' sowie einige HTML5-Funktionen. Erkennbar ist dies z.B. an dem Attribut "`required"', welches dem Browser signalisiert, dass dieses Formularelement ausgef�llt werden muss. Versucht der Benutzer das Formular abzuschicken erh�lt er einen Hinweis �ber das leere Element und wird aufgefordert es auszuf�llen - vorausgesetzt der Browser ist HTML5-kompatibel. Das Formular wird durch eine �bliche Tabelle strukturiert (table-tag).

%CODE 
/*...*/

			<tr>
				<td>Name: *</td>
				<td><input id="name" placeholder="Bitte Name angeben" type="text" tabindex="1" required autofocus></td>
			</tr>
/*...*/
%CODE Ende

F�r den Fall einer fehlenden HTML5-Unterst�tzung wurde eine Fallback-L�sung mittels JavaScript implementiert. Die Pr�fung auf eine HTML5-Unterst�tzung wurde wie folgt umgesetzt und wird bei jedem Abschicken des Formulars durchgef�hrt.

%CODE
/*...*/
	var inputs = document.createElement('input');
	var supports = {};
	supports.required  = 'required' in inputs;

	if(!supports.required) {
	/*...siehe unten...*/
	}
/*...*/
%CODE Ende

Trifft die Pr�fung zu und HTML5 wird nicht unterst�tzt wird jedes auszuf�llende Formularelement auf dessen Inhalt gepr�ft. Bei einem Fehler wird die Formular�bermittlung abgebrochen und der Fokus auf das betroffene Element gesetzt.

%CODE

/*...*/
		 name = document.getElementById("name").value;
/*...*/
		 if (name == "")
		 {
		 document.getElementById("name").select();
		 document.getElementById("name").focus();
		 return false;
		 }
/*...*/

%CODE ENDE

Da HTML5 auch die Plausibilit�t von E-Mail-Adressen pr�ft wurde auch dies im Fallback ber�cksichtigt. Hierf�r wurde ein Regul�rer Ausdruck verwendet:

%CODE
/*...*/
		 if (!checkEmail(email))
		 {
		 alert('E-Mail ist ung�ltig.');
		 return false;
		 }
/*...*/

function checkEmail(inputvalue) {
	var pattern = /^([a-zA-Z0-9_\.\-])+\@(([a-zA-Z0-9\-])+\.)+([a-zA-Z0-9]{2,4})+$/;
	return pattern.test(inputvalue);
}
/*...*/

%CODE ENDE
%%\section{Beispielseite: Informationen zur Unterkunft}

\textsc{Mathias Vongerichten bzw. Stefan Waidele}

\subsection{�berschrift}

Erl�uterungen zur redaktionellen, gestalterischen und technischen Umsetzung.


%%\section{Beispielseite: Informationen �ber die Busse}

\textsc{Matthias Vongerichten bzw. Stefan Waidele}

\subsection{�berschrift}

Erl�uterungen zur redaktionellen, gestalterischen und technischen Umsetzung.


%%\section{Beispielseite: Unternehmensportrait}

\textsc{Mathias Vongerichten bzw. Stefan Waidele}

\subsection{�berschrift}

Erl�uterungen zur redaktionellen, gestalterischen und technischen Umsetzung.


%%\input{beispielimpressum.tex}
\section{Fazit}

\textsc{Matthias Vongerichten}

Das redaktionelle Konzept beschreibt die vorgesehenen Inhalte je Seite in Kurzform. Au�erdem wird hierbei auch eine Richtlinie zur Aktualisierungsh�ufigkeit je Seite vorgegeben (unterst�tzt SEO). Somit wird erreicht, dass der redaktionelle Verantwortliche eine kompakte Anleitung zur Pflege der Website erh�lt.

Beim Navigationskonzept wurden die wichtigsten Navigationselemente beachtet. Hervorzuheben sind hier die "`Breadcrumb"'-Navigation und die Implementierung von mod\_rewrite (Abbildung ~\ref{abb:htaccess}, Seite ~\pageref{abb:htaccess}).

\textsc{Stefan Waidele}

...

\subsection{Evt. Gesamtfazit}

Die Website erf�llt die gestellten Anforderungen vollst�ndig. So bescheinigt uns woorank.com (SEO-Analyse-Tool), dass z.B. nicht-funktionalen Anforderungen wie die Performance der Site sehr gut erf�llt werden.

\begin{figure}[h]
\includegraphics[width=\hsize]{scr_woorank_performance.png}
\caption{Screenshot: Performance}
\label{scr:performance}
\end{figure}
\begin{figure}[h]
\includegraphics[width=\hsize]{scr_woorank_mobile_performance.png}
\caption{Screenshot: Mobile Ger�te}
\label{scr:mobile}
\end{figure}

%screenshot scr_woorank_performance.png oder scr_woorank_mobile performance.png oder beides

Ein n�chster m�glicher Schritt ist die Verwendung eines Content-Management-Systems (CMS). Dies erm�glicht das einfache Bearbeiten der Seiteninhalte mithilfe eines WYSIWYG-Editors (keine Code-Kenntnisse notwendig) und das Aufrechterhalten eines konsistenten Designs. Die Navigation kann auf mehrere Ebenen und Elemente verteilt werden und aktualisiert sich selbst�ndig. Der Administrationsaufwand wird somit stark reduziert und trotzdem f�llt es erstaunlich leicht die Funktionen der Website durch vorgefertigte Module zu erweitern. So ist es ebenfalls m�glich einen Weblog zu betreiben welcher mit Facebook und anderen Social-Media-Diensten integriert werden kann. Urspr�nglich als Weblog-System zu gro�er Bekanntheit gelangt, kann Wordpress heute als einfaches und komfortables CMS genannt werden. F�r komplexere Anwendungsf�lle eignet sich Drupal oder Typo3.

\end{onehalfspacing}

%% Bibliographie unter Verwendung von dinnat %%%%%%%%%%%%%%%%%%%%%%%%%%
%\setbibpreamble{Präambel}		% Text vor dem Verzeichnis
%\bibliographystyle{natdin}
%\bibliographystyle{plaindin}
%\bibliography{Literatur}	% Sie benötigen eine *.bib-Datei

%% Bibliographie unter Verwendung von BibLaTeX %%%%%%%%%%%%%%%%%%%%%%%%%%
\clearpage
\flushleft
\setlength{\bibitemsep}{1em}
\printbibliography

\clearpage

\pagestyle{empty} 
\thispagestyle{empty}

\textsc{\Large Eidesstattliche Erkl"arung}
\vspace*{4cm}

Ich versichere, dass ich das beiliegende Assignment selbstst"andig verfasst, keine anderen als die angegebenen Quellen und Hilfsmittel benutzt sowie alle w"ortlich oder sinngem"a"s "ubernommenen Stellen in der Arbeit gekennzeichnet habe. 
\vspace{3cm}

\begin{tabbing}
Wherever, \today \hspace{1em} \=\rule[-0.8ex]{6.5cm}{0.75pt}\\
                \>(Matthias Vongerichten) 
\end{tabbing}
\vspace{3cm}

\begin{tabbing}
Neuenburg am Rhein, \today \hspace{1em} \=\rule[-0.8ex]{6.5cm}{0.75pt}\\
                \>(Stefan Waidele) 
\end{tabbing}


\end{document}
